\section*{Devoir maison }

\begin{Exercise}[title={Un peu de binaire}]
\Question Donner la représentation binaire sur un octet des nombres suivants, justifiez vos réponses par un calcul :
\begin{itemize}
    \item 2 :
    \item 156 :
    \item 89 :
\end{itemize}
\Question Donner la représentation décimale des nombres suivants (ici le bit de poids fort est à gauche), justifiez vos réponses par un calcul :
\begin{itemize}
    \item 1000001 :
    \item 101 :
    \item 111111110 :
\end{itemize}
\end{Exercise}

\begin{Exercise}[title={Conversion binaire}]
On modélise la représentation binaire d'un entier non signé par une liste d'entiers dont les éléments sont 0 ou 1.\\
Par exemple, le tableau [1, 0, 1, 0, 0, 1, 1] représente
l'écriture binaire de l'entier dont l'écriture décimale est
$2^6 + 2^4 + 2^1 + 2^0 = 83$.\\
À l'aide d'un parcours par index du tableau, écrire la fonction convertir répondant aux spécifications suivantes :
\begin{lstlisting}[language=Python]
def convertir(tab):
"""
tab est un tableau d'entiers, dont les éléments sont 0 ou 1,
et représentant un entier écrit en binaire.
Renvoie l'écriture décimale de l'entier positif dont la
représentation binaire est donnée par le tableau tab
"""
.....................................................................
.....................................................................
.....................................................................
.....................................................................
.....................................................................
.....................................................................
.....................................................................
.....................................................................
.....................................................................
.....................................................................
.....................................................................
.....................................................................
.....................................................................
.....................................................................
\end{lstlisting}

Exemples d'appels :
\begin{itemize}
    \item convertir([1, 0, 1, 0, 0, 1, 1]) retourne 83
    \item convertir([1, 0, 0, 0, 0, 0, 1, 0]) retourne 130
\end{itemize}
\end{Exercise}

